\section{From Edges to Nodes}
\label{sec:edges-to-nodes}

\subsection{Balance Equation}
\label{subsec:balance-equation}

With the graph convention from Section~\ref{sec:edges-and-nodes}, a node view
becomes available when an edge price is interpreted as a ratio of latent node
values.  Let $\nu_u \in \Rpos$ denote the latent value of node
$u \in \mathcal{V}$.  For edge $e=(u,v) \in \mathcal{E}$, the idealized edge
price is
\[
  q_e = \frac{\nu_u}{\nu_v}
  \in \Rpos.
\]
In log coordinates,
\[
  \ell_e := \log q_e = \xi_u - \xi_v,
  \qquad
  \ell_e \in \R,
  \qquad
  \xi_u := \log \nu_u \in \R.
\]
Therefore the signed edge return satisfies
\[
  r_e = \Delta \ell_e = y_u - y_v,
  \qquad
  r_e \in \R,
  \qquad
  y_u := \Delta \xi_u \in \R.
\]
This is the basic balance equation.  The edge forecast is a relation; the node
state is a potential whose differences explain that relation.

\subsection{Best-Fit Node Recovery}
\label{subsec:best-fit-node-recovery}

In practice the active edge set may be incomplete, noisy, or internally
inconsistent.  The node values are then recovered as the best global fit to the
observed or predicted edge returns.

For one forecast step and one active scalar observer, let
$\widehat{r}_e \in \R$ be the observed edge return and let
$\gamma_e \in \Rpos$ be its confidence weight.  Recover the node vector
$\vect{y} \in \R^N$ by
\[
  \min_{\vect{y} \in \R^N}
  \sum_{e=(u,v) \in \mathcal{E}}
  \gamma_e\left(\widehat{r}_e - (y_u - y_v)\right)^2
  \qquad\text{subject to}\qquad
  y_g = 0.
\]
The gauge constraint removes the additive ambiguity: only differences of node
values matter.

\subsection{Incidence-Matrix Form}
\label{subsec:incidence-matrix-form}

Let $\mat{A} \in \R^{L \times N}$ be the directed incidence matrix of
$\mathcal{G}$.  For edge $e=(u,v)$, the corresponding incidence entries are
\[
  A_{e\,u}=1,\qquad A_{e\,v}=-1,
\]
and zeros elsewhere.  Let $\vect{\widehat{r}} \in \R^L$ collect the edge
observations and let
\[
  \mat{\Gamma} = \diag(\gamma_1,\ldots,\gamma_L)
  \in \R^{L \times L}.
\]
Then node recovery is the weighted least-squares problem
\[
  \min_{\vect{y} \in \R^N}
  \lVert \mat{\Gamma}^{1/2}(\vect{\widehat{r}} - \mat{A}\vect{y}) \rVert_2^2,
  \qquad y_g = 0.
\]
Its normal equations are
\[
  \mat{A}^{\top}\mat{\Gamma}\mat{A}\vect{y}
  =
  \mat{A}^{\top}\mat{\Gamma}\vect{\widehat{r}}.
\]
The matrix
\[
  \mat{L}_{\gamma} := \mat{A}^{\top}\mat{\Gamma}\mat{A}
  \in \R^{N \times N}
\]
is the weighted graph Laplacian of the active edge universe in the usual
incidence-matrix form.

\subsection{Covariance-Aware Recovery}
\label{subsec:covariance-aware-recovery}

If a full edge covariance estimate $\mat{C}_r \in \R^{L \times L}$ is available
for the chosen observer, then the natural generalization is
\[
  \min_{\vect{y} \in \R^N}\ (
  \vect{\widehat{r}} - \mat{A}\vect{y}
  )^{\top}
  \mat{C}_r^{-1}
  (
  \vect{\widehat{r}} - \mat{A}\vect{y}
  ),
  \qquad y_g = 0.
\]
This is the Mahalanobis-weighted version of the same balance equation.  It does
not introduce a new idea; it only measures mismatch using the edge uncertainty
structure instead of independent scalar weights.

\subsection{Node Covariance and Geometry}
\label{subsec:node-covariance-and-geometry}

When the covariance-aware system is well posed, the recovered node covariance is
approximated by
\[
  \mat{C}_y \approx (\mat{A}^{\top}\mat{C}_r^{-1}\mat{A})^{-1},
  \qquad
  \mat{C}_y \in \R^{N \times N},
\]
after gauge fixing or the equivalent constrained solve.  This gives a natural
local geometry on node states.  For two node configurations
$\vect{y},\vect{y}' \in \R^N$, the Mahalanobis distance is
\[
  d(\vect{y},\vect{y}')^2
  =
  (\vect{y} - \vect{y}')^{\top}
  \mat{C}_y^{-1}
  (\vect{y} - \vect{y}')
  \in \Rnonneg.
\]
This distance measures separation in units of node uncertainty, not in raw
coordinate scale.

\subsection{Modes}
\label{subsec:modes}

If a covariance matrix is available on the recovered node space, diagonalize it
as
\[
  \mat{C}_y = \mat{V}\Lambda\mat{V}^{\top},
  \qquad
  \mat{V} \in \R^{N \times N},
  \qquad
  \Lambda \in \R^{N \times N}.
\]
The eigenvectors are the principal modes of collective node motion, and the
eigenvalues measure how much variance each mode carries.  The formal dimension
remains $N$; keeping only a few dominant modes is an optional compression step,
not a mandatory reduction.
